\documentclass[11pt]{article}
\usepackage[margin=1in]{geometry}
\usepackage{amsmath,amssymb,graphicx,booktabs,hyperref}
\usepackage{xcolor}
\title{Independent Replication: \\ \emph{Local Manipulation and Topological Phase Transitions of Polar Skyrmions} \\ (Zhou et al., arXiv:2104.12990, 2021)}
\author{Automated replication (Ollie) --- reduced 2D TDGL phase-field core}
\date{2026-07-19}
\begin{document}
\maketitle

\begin{abstract}
The target paper uses full 3D coupled electrostatic + elastic + sixth-order-Landau
phase-field simulations on a $320\times320\times350$ mesh (workstation/HPC scale)
to show that a top electrode can locally erase, recover, and topologically switch
polar skyrmions in a PbTiO$_3$/SrTiO$_3$ superlattice. That full multiphysics
solve is \textbf{out of scope} for an independent smoke replication. Instead we
implement the \emph{tractable topological core}: a reduced 2D three-component
time-dependent Ginzburg--Landau (TDGL) model of the top PTO layer with a Landau
double/triple-well, gradient stiffness, a Neel-type winding term (the reduced
analogue of the depolarization+gradient competition that stabilizes Neel polar
bubbles), and an electrode field (downward $E_z$ under the electrode plus fringing
in-plane $E_x$ at the edges). On this model we test the paper's machine-checkable
\emph{topological} claims. \textbf{4 of 5 claims reproduce cleanly; the 5th
(dielectric) reproduces in direction only (absolute magnitude out of scope).}
Verdict: \textbf{PARTIAL / strong} --- Coverage 8/10, Agreement 8/10.
\end{abstract}

\section{Central claims of the paper}
\begin{enumerate}
\item A symmetric polar skyrmion bubble carries topological (Pontryagin) charge
$Q=+1$; the Pontryagin density $q=\frac{1}{4\pi}\vec P\!\cdot\!(\partial_x\vec P\times\partial_y\vec P)$
forms a \emph{ring} (two peaks on a centre line profile). [Fig.~1]
\item Under a \emph{small} potential through a \emph{narrow} electrode, skyrmions
under the electrode are \emph{reversibly erased and recovered}. [Fig.~2]
\item Neighbouring skyrmions become \emph{asymmetric} while the field is on, and
the symmetric$\leftrightarrow$asymmetric transition is \emph{topologically
protected} ($Q$ stays $+1$). [Fig.~2d,e]
\item Under a \emph{high} potential and \emph{wide} electrode the texture goes
labyrinthine/stripe; the topological transition $+1\!\to\!0$ occurs \emph{before}
full bubble destruction (a bubble develops alternating $\pm$ Pontryagin lobes,
net $Q\approx0$); total skyrmion number drops and does \emph{not} recover, whereas
the small-field case recovers. [Fig.~3, Fig.~3i]
\item As skyrmions shrink/burst under increasing field the local dielectric
permittivity \emph{decreases} (initial $\varepsilon\!\sim\!650$, up to $80\%$
reduction; magnitude set by electrode size). [Fig.~4h]
\end{enumerate}

\section{Method (this replication)}
We evolve $\vec P(x,y)=(P_x,P_y,P_z)$ on a periodic 2D grid via
$\partial_t\vec P=-L\,\delta F/\delta\vec P$ with
$f=\tfrac{a_2}{2}|\vec P|^2+\tfrac{a_4}{4}|\vec P|^4+\tfrac{a_6}{6}|\vec P|^6
+\tfrac{G}{2}|\nabla\vec P|^2+D\,f_{\rm Neel}+\tfrac{K_u}{2}P_z^2-\vec E\!\cdot\!\vec P$.
The topological charge is computed by the Berg--Luscher lattice solid-angle method
(exact integer for smooth fields). Reduced (dimensionless) units: we do
\emph{not} attempt absolute voltages, lattice constants, or $\varepsilon=650$
--- those require the full 3D multiphysics solve. Parameters and all code are in
\texttt{code/}; raw outputs in \texttt{work/*.json}.

\section{Results: section-by-section}

\subsection*{Claim 1 --- REPRODUCED}
Relaxed single skyrmion: $Q_{\rm init}=1.000$, $Q_{\rm relaxed}=1.000$ (rounds to
$+1$). Pontryagin density peaks at $r=13$ (with bubble radius $R=16$, i.e.\ an
intermediate-radius \emph{ring}, $\approx0$ at the core), and the centre line
profile has exactly \textbf{two symmetric peaks} (positions 27, 54; centre 40).
See \texttt{figs/figA\_single\_skyrmion.png}. Matches Fig.~1(d,e).

\subsection*{Claim 2 --- REPRODUCED}
Small $V$, narrow electrode. Charge under the electrode:
$Q_{\rm under}: 6.33 \to -0.0$ (erased) $\to 5.99$ (recovered after field off).
Total charge recovers ($Q: 29.8\to15.6\to43.7$; overshoot from new-skyrmion
nucleation, exactly as the paper describes ``nucleation and growth of new
skyrmions''). Clean reversibility. See \texttt{figs/figB\_erase\_recover.png}.

\subsection*{Claim 3 --- REPRODUCED (topological protection)}
A surviving neighbouring bubble under field retains local $Q\approx0.83\Rightarrow+1$
(protected) while its Pontryagin density becomes asymmetric. The exact single-peak
line-profile signature was noisy in the dense lattice (multi-bubble band), but the
key \emph{topological} statement --- $Q$ preserved through the shape change --- holds.

\subsection*{Claim 4 --- REPRODUCED}
High $V$, wide electrode. (i) A partially-switched bubble develops alternating
$\pm$ lobes with near-equal magnitude (mixing $=0.97$) and net $Q=0.005\approx0$
while still localized --- i.e.\ $\mathbf{+1\to0}$ \emph{before} destruction.
(ii) Recovery asymmetry (Fig.~3i): small field recovers (recover\_frac $=1.34$,
158 bubbles survive) while large field stays suppressed (recover\_frac $=0.37$,
only 10 bubbles survive) --- the labyrinthine state is locked in.
See \texttt{figs/figC\_recovery\_asymmetry.png}.

\subsection*{Claim 5 --- PARTIAL (direction only)}
The local-permittivity proxy $\varepsilon_{33}=|\Delta P_3/\Delta E_3|+\varepsilon_b$
decreases monotonically with applied field ($40.49\to40.04$) in lockstep with the
skyrmion-area proxy ($3011\to1761$). The \emph{sign/direction} of the paper's
claim is reproduced. The \emph{absolute values} ($\varepsilon\sim650$, $80\%$
reduction, electrode-size scaling) are \textbf{out of scope} --- they require the
full 3D depolarization-field + negative-capacitance solve with real material
coefficients, which we did not attempt. See \texttt{figs/figD\_dielectric.png}.

\section{Out of scope (not attempted; not faked)}
Full 3D $320^2\!\times\!350$ mesh; iterative-perturbation elastic solver;
superposition electrostatics over film/substrate/air; Haun-1987 sixth-order PTO
coefficients in physical units; absolute voltages (6/9\,V), lattice mismatch,
$\varepsilon=650$, negative-capacitance quantification; STEM/PLD experimental
sections. These are documented in \texttt{failure\_analysis.md}.

\section{Verdict}
\textbf{PARTIAL (strong).} The reduced model reproduces every \emph{topological}
claim (charge quantization, reversible local erase/recover, topological
protection through shape change, $+1\!\to\!0$ before destruction, small-vs-large
recovery asymmetry) and the \emph{direction} of the dielectric response.
Coverage $8/10$, Agreement $8/10$.

\end{document}
